% ===== §9 =====
\section{The Insufficiency of the Functional Case}
\label{sec:functionlimit}

\noindent
The functional case sets out what relational understanding does, the trust it generates, the shared value it recognizes, the conflict it eases, the quality of a shared life it sustains. A case of this form defends a practice by its effects, and a defense by effects carries a weakness that must be faced directly, since it bears on the whole standing of the account. The weakness is that the effects by which the practice is defended can, in some measure and for some time, be produced by something that is not the practice at all, and that a defense by effects alone has no way to tell the two apart.

Consider what an efficient possession of another can counterfeit. A partner who studies the other to secure a hold upon the other's world and not to enter it, who learns the other's needs and fears in order to manage them, who anticipates and arranges so that the other is kept content and compliant, may produce, for a time and by the measures the functional case employs, much of what relational understanding produces. Such a partner may generate a felt trust, since the other, well managed, finds their expectations met. Such a partner may produce an appearance of shared value, having learned to present as common what secures the other's attachment. Such a partner may ease conflict, having learned which crossings to avoid and which concessions to trade. Such a partner may sustain a certain satisfaction in the bond, meeting the other's needs with the efficiency that a studied management affords. By the functions alone, this bond and a bond of genuine relational understanding may for a considerable time be difficult to tell apart, and on some measures the managed bond may score higher, since management is directed at the measures in a way that understanding is not.

This is the weakness of the functional case, and it is not a marginal one. It is the same weakness that attends every defense of a practice by its effects, that the effects may be counterfeited by a simulacrum directed at producing them, and that the defense, having appealed only to the effects, cannot distinguish the practice from its simulacrum. In the domain of intimacy the simulacrum has a name in the wider framework within which this paper stands. It is a relational extraction, a taking from the other under the form of a giving, and its refinement is precisely its danger, since the more skillfully the functions of understanding are counterfeited the harder the extraction is to detect and the longer it may run.

The point is not that the functional case is false. The functions are real, and relational understanding does produce them. The point is that the functional case is not sufficient, because it cannot, from within itself, distinguish the understanding that produces these functions as the overflow of a genuine crossing from the management that produces their appearance as the instrument of a hold. To distinguish these is to appeal to something other than the functions, to what the practice is and not only to what it yields, and this is a normative and not a functional distinction. The remainder of this part draws it.

It is worth marking why the distinction cannot be recovered by a longer or a more careful functional analysis. One might propose that the counterfeit will eventually fail on the measures, that management, unlike understanding, cannot sustain the functions indefinitely, so that a sufficiently long functional accounting would tell them apart after all. There is something in this, and the next parts will give it its due, since the temporal structure of intimacy does bear on what can be sustained. Yet it does not rescue the functional case, for two reasons. The counterfeit may run for a very long time, across years in which real harm is done, before any functional shortfall appears, so that a distinction available only in the eventual accounting comes too late to guide a life. And the deeper reason is that even where the counterfeit does eventually fail on the measures, what it lacked was never merely a quantity of function. It lacked something the measures do not capture, for want of which we judge it a wrong done to the other and not merely a bond that underperformed. That something is what the normative case must name.
